\documentclass[a4paper,12pt]{article}

\usepackage[T2A]{fontenc}
\usepackage[utf8]{inputenc}
\usepackage[russian]{babel}

\usepackage{amsmath, amssymb, amsthm, mathtools}
\usepackage{helvet}
\renewcommand{\familydefault}{\sfdefault}
\usepackage{icomma}
\usepackage{geometry}
\usepackage{microtype}
\usepackage{tikz}
\usepackage{tkz-euclide}
\usepackage{pgfplots}
\pgfplotsset{compat=1.18}
\usepackage{tabularx, booktabs, longtable}
\usepackage{enumitem}
\usepackage{hyperref}
\usepackage{parskip}
\usepackage{fancyhdr}
\usepackage{setspace}
\usepackage{xcolor}

\geometry{left=18mm,right=18mm,top=18mm,bottom=20mm}
\onehalfspacing
\sloppy

\definecolor{accent}{HTML}{008080}
\definecolor{darkaccent}{HTML}{004D4D}
\definecolor{textgray}{HTML}{333333}
\definecolor{softgray}{HTML}{F4F6F6}

\hypersetup{
    colorlinks=true,
    linkcolor=darkaccent,
    urlcolor=darkaccent
}

\pagestyle{fancy}
\fancyhf{}
\lhead{\textcolor{textgray}{Чек-лист: текстовые задачи}}
\rhead{\textcolor{textgray}{ОГЭ}}
\cfoot{\thepage}
\renewcommand{\headrulewidth}{0.4pt}
\renewcommand{\headrule}{\hbox to\headwidth{\color{accent}\leaders\hrule height \headrulewidth\hfill}}

\setlist[itemize]{leftmargin=*, itemsep=2pt, topsep=3pt}
\setlist[enumerate]{leftmargin=*, itemsep=4pt, topsep=4pt}

\newcommand{\important}[1]{\textcolor{darkaccent}{\textbf{#1}}}
\newcommand{\kmh}{\text{км/ч}}
\newcommand{\km}{\text{км}}
\newcommand{\ch}{\text{ч}}
\newcommand{\minut}{\text{мин}}
\newcommand{\ssec}{\text{с}}
\newcommand{\m}{\text{м}}
\newcommand{\lpm}{\text{л/мин}}

\newcommand{\trainpole}{
\begin{center}
\begin{tikzpicture}[scale=0.95]
    \draw[line width=0.8pt] (0,0) -- (9,0);
    \draw[line width=1pt, color=darkaccent] (1,0.2) rectangle (4.3,0.9);
    \foreach \x in {1.5,2.5,3.5} {
        \fill[accent] (\x,0.2) circle (2pt);
    }
    \draw[line width=1pt] (6.2,0) -- (6.2,1.4);
    \fill[darkaccent] (6.2,1.4) circle (2pt);
    \draw[->, line width=1pt, color=darkaccent] (4.6,0.55) -- (5.8,0.55);
    \node[above] at (2.65,0.9) {поезд};
    \node[above] at (6.2,1.4) {столб};
    \node[below] at (2.65,0.05) {$S_1$};
    \node[below] at (6.2,0) {$S_2=0$};
\end{tikzpicture}
\end{center}
}

\newcommand{\twotrains}{
\begin{center}
\begin{tikzpicture}[scale=0.95]
    \draw[line width=0.8pt] (0,0) -- (9,0);
    \draw[line width=1pt, color=darkaccent] (0.8,0.25) rectangle (3.2,0.9);
    \draw[line width=1pt, color=darkaccent] (6.2,0.25) rectangle (8.5,0.9);
    \draw[->, line width=1pt, color=accent] (3.4,0.58) -- (4.7,0.58);
    \draw[->, line width=1pt, color=accent] (6.0,0.58) -- (4.7,0.58);
    \node[above] at (2,0.9) {$S_1,\ v_1$};
    \node[above] at (7.35,0.9) {$S_2,\ v_2$};
    \node[below] at (4.65,0) {навстречу: скорости складываются};
\end{tikzpicture}
\end{center}
}

\newcommand{\bargeovertake}{
\begin{center}
\begin{tikzpicture}[scale=0.92]
    \draw[line width=0.8pt, color=accent] (0,0) -- (10,0);
    \draw[line width=1pt, color=darkaccent] (0.8,0.35) rectangle (2.4,0.85);
    \draw[line width=1pt, color=darkaccent] (4.6,0.35) rectangle (7.0,0.85);
    \draw[<->, color=accent] (2.4,1.15) -- (4.6,1.15);
    \node[above] at (3.5,1.15) {$S_{\text{дог}}$};
    \node[below] at (1.6,0.35) {$S_1$};
    \node[below] at (5.8,0.35) {$S_2$};
    \draw[->, line width=1pt, color=darkaccent] (2.6,0.6) -- (3.8,0.6);
    \draw[->, line width=1pt, color=darkaccent] (7.2,0.6) -- (8.1,0.6);
    \node[below] at (5,-0.2) {догон и перегон учитываются как дополнительный относительный путь};
\end{tikzpicture}
\end{center}
}

\begin{document}

\begin{titlepage}
\thispagestyle{empty}
\begin{center}
\vspace*{25mm}

{\Large \textcolor{accent}{\textbf{Чек-лист}}}

\vspace{8mm}

{\Huge \textbf{Текстовые задачи}}

\vspace{4mm}

{\Large средняя скорость, протяжённые объекты, производительность}

\vspace{22mm}

{\large Лёвин Артём Александрович\\[1mm]
эксклюзивно для Софии Кальней}

\vspace{12mm}

{\large \today}

\vfill

\begin{tikzpicture}
    \draw[color=accent, line width=1.2pt]
    (0,0) .. controls (2,0.8) and (4,-0.8) .. (6,0)
          .. controls (8,0.8) and (10,-0.8) .. (12,0);
    \draw[color=darkaccent, line width=0.5pt]
    (1,-0.35) .. controls (3,0.25) and (5,-0.95) .. (7,-0.35)
          .. controls (9,0.25) and (10.5,-0.75) .. (11.5,-0.35);
\end{tikzpicture}

\vspace{8mm}
\end{center}
\end{titlepage}

\tableofcontents
\newpage

\section{Средняя скорость}

\subsection{Главная формула}

Средняя скорость считается не как обычное среднее арифметическое скоростей, а как отношение всего пути ко всему времени:
\[
v_{\text{ср}}=\frac{S_{\text{общ}}}{t_{\text{общ}}}.
\]

Если движение состоит из нескольких участков, то:
\[
v_{\text{ср}}=\frac{S_1+S_2+\cdots+S_n}{t_1+t_2+\cdots+t_n}.
\]

\important{Ключевая мысль:} в числителе всегда общий путь, в знаменателе всегда общее время.

\subsection{Случай 1. Половина времени}

Если первую половину времени объект двигался со скоростью $84\ \kmh$, а вторую половину времени --- со скоростью $56\ \kmh$, то время удобно обозначить так:
\[
\frac{t}{2} \quad \text{и} \quad \frac{t}{2}.
\]

\begin{center}
\renewcommand{\arraystretch}{1.35}
\begin{tabularx}{0.82\linewidth}{X c c c}
\toprule
 & $v$ & $t$ & $S$ \\
\midrule
первый участок & $84$ & $\dfrac{t}{2}$ & $84\cdot\dfrac{t}{2}=42t$ \\
второй участок & $56$ & $\dfrac{t}{2}$ & $56\cdot\dfrac{t}{2}=28t$ \\
\bottomrule
\end{tabularx}
\end{center}

Средняя скорость:
\[
v_{\text{ср}}=\frac{42t+28t}{\frac{t}{2}+\frac{t}{2}}.
\]
\[
v_{\text{ср}}=\frac{70t}{t}=70.
\]

\[
\boxed{70\ \kmh}
\]

\important{Важный вывод:} если равны именно времена движения, то средняя скорость равна обычному среднему арифметическому скоростей:
\[
v_{\text{ср}}=\frac{84+56}{2}=70.
\]

\subsection{Случай 2. Половина пути}

Если первую половину пути объект ехал со скоростью $56\ \kmh$, а вторую половину пути --- со скоростью $84\ \kmh$, то путь удобно обозначить так:
\[
\frac{S}{2} \quad \text{и} \quad \frac{S}{2}.
\]

\begin{center}
\renewcommand{\arraystretch}{1.35}
\begin{tabularx}{0.82\linewidth}{X c c c}
\toprule
 & $v$ & $t$ & $S$ \\
\midrule
первый участок & $56$ & $\dfrac{S}{112}$ & $\dfrac{S}{2}$ \\
второй участок & $84$ & $\dfrac{S}{168}$ & $\dfrac{S}{2}$ \\
\bottomrule
\end{tabularx}
\end{center}

Средняя скорость:
\[
v_{\text{ср}}=
\frac{\frac{S}{2}+\frac{S}{2}}{\frac{S}{112}+\frac{S}{168}}.
\]

Так как
\[
\frac{S}{2}+\frac{S}{2}=S,
\]
то:
\[
v_{\text{ср}}=\frac{S}{S\left(\frac{1}{112}+\frac{1}{168}\right)}.
\]

\[
\frac{1}{112}+\frac{1}{168}
=
\frac{3}{336}+\frac{2}{336}
=
\frac{5}{336}.
\]

\[
v_{\text{ср}}=\frac{S}{S\cdot \frac{5}{336}}=\frac{336}{5}=67{,}2.
\]

\[
\boxed{67{,}2\ \kmh}
\]

\important{Важный вывод:} если равны пути, средняя скорость не равна обычному среднему арифметическому скоростей.

\section{Протяжённые объекты: поезда, платформы, колонны}

\subsection{Почему поезд нельзя считать точкой}

Автомобиль или пешеход в большинстве задач можно считать точкой. Поезд, баржу, колонну или состав нельзя считать точкой, потому что у них есть длина.

Если поезд проходит мимо столба, он должен полностью пройти свою длину.

\trainpole

Если поезд проходит мимо платформы, он должен пройти:
\[
\text{длина поезда}+\text{длина платформы}.
\]

\subsection{Базовая формула}

Для времени прохождения одного протяжённого объекта мимо другого используется формула:
\[
t=\frac{S_1+S_2}{v_{\text{отн}}},
\]
где:
\[
S_1,\ S_2 \text{ --- длины объектов},
\]
\[
v_{\text{отн}} \text{ --- относительная скорость}.
\]

Если объекты движутся навстречу друг другу, скорости складываются:
\[
v_{\text{отн}}=v_1+v_2.
\]

Если один объект догоняет другой, скорости вычитаются:
\[
v_{\text{отн}}=v_1-v_2.
\]

\important{Если объект неподвижен, его скорость равна нулю.}

\important{Если объект практически не имеет длины по сравнению с поездом, его длину можно считать равной нулю.}

\subsection{Пример 1. Поезд проходит мимо столба}

Поезд, скорость которого равна $60\ \kmh$, проходит мимо столба за $30$ секунд. Найдите длину поезда в метрах.

Переводим время в часы:
\[
30\ \ssec=\frac{30}{3600}=\frac{1}{120}\ \ch.
\]

Пусть $x$ --- длина поезда в километрах. Длина столба в направлении движения считается равной нулю:
\[
S_2=0.
\]

\[
t=\frac{S_1+S_2}{v_1+v_2}.
\]

\[
\frac{1}{120}=\frac{x+0}{60+0}.
\]

\[
x=\frac{60}{120}=0{,}5\ \km.
\]

Переводим в метры:
\[
0{,}5\ \km=500\ \m.
\]

\[
\boxed{500\ \m}
\]

\subsection{Пример 2. Поезд и пешеход движутся навстречу}

Поезд и пешеход движутся навстречу друг другу. Скорость поезда равна $60\ \kmh$, скорость пешехода равна $3\ \kmh$. Поезд проходит мимо пешехода за $30$ секунд. Найдите длину поезда.

Длина пешехода в направлении движения считается равной нулю:
\[
S_2=0.
\]

Так как движение навстречу, скорости складываются:
\[
v_{\text{отн}}=60+3=63\ \kmh.
\]

\[
30\ \ssec=\frac{30}{3600}=\frac{1}{120}\ \ch.
\]

\[
\frac{1}{120}=\frac{x}{63}.
\]

\[
x=\frac{63}{120}=0{,}525\ \km.
\]

\[
0{,}525\ \km=525\ \m.
\]

\[
\boxed{525\ \m}
\]

\subsection{Пример 3. Два поезда движутся навстречу}

Если два поезда движутся навстречу друг другу, то для полного прохождения друг мимо друга они вместе должны покрыть сумму своих длин.

\twotrains

Пусть длины поездов равны $400\ \m$ и $500\ \m$, а скорости равны $72\ \kmh$ и $90\ \kmh$. Найдите время прохождения поездов друг мимо друга.

Переводим длины в километры:
\[
400\ \m=0{,}4\ \km, \qquad 500\ \m=0{,}5\ \km.
\]

Сумма длин:
\[
0{,}4+0{,}5=0{,}9\ \km.
\]

Относительная скорость:
\[
72+90=162\ \kmh.
\]

\[
t=\frac{0{,}9}{162}=\frac{1}{180}\ \ch.
\]

Переводим в секунды:
\[
\frac{1}{180}\cdot 3600=20.
\]

\[
\boxed{20\ \ssec}
\]

\section{Догон и перегон протяжённых объектов}

\subsection{Расширенная формула}

Если один протяжённый объект догоняет другой, а затем перегоняет его на некоторое расстояние, нужно учитывать не только длины объектов, но и дополнительные расстояния.

\bargeovertake

Расширенная формула:
\[
t=
\frac{S_1+S_2+S_{\text{дог}}+S_{\text{пер}}}{v_{\text{быстр}}-v_{\text{медл}}}.
\]

Здесь:
\[
S_1,\ S_2 \text{ --- длины объектов},
\]
\[
S_{\text{дог}} \text{ --- начальное отставание},
\]
\[
S_{\text{пер}} \text{ --- расстояние, на которое нужно перегнать},
\]
\[
v_{\text{быстр}}-v_{\text{медл}} \text{ --- разность скоростей}.
\]

\subsection{Пример}

Короткая баржа длиной $40\ \m$ сначала отстаёт от длинной баржи длиной $60\ \m$ на $200\ \m$. Затем короткая баржа догоняет длинную и перегоняет её на $300\ \m$. На это уходит $18$ минут. На сколько $\kmh$ скорость короткой баржи больше скорости длинной?

Переводим все расстояния в километры:
\[
40\ \m=0{,}04\ \km,
\]
\[
60\ \m=0{,}06\ \km,
\]
\[
200\ \m=0{,}2\ \km,
\]
\[
300\ \m=0{,}3\ \km.
\]

Общий относительный путь:
\[
0{,}04+0{,}06+0{,}2+0{,}3=0{,}6\ \km.
\]

Переводим время:
\[
18\ \minut=\frac{18}{60}=0{,}3\ \ch.
\]

Пусть
\[
\Delta v=v_{\text{короткой}}-v_{\text{длинной}}.
\]

Тогда:
\[
0{,}3=\frac{0{,}6}{\Delta v}.
\]

\[
\Delta v=\frac{0{,}6}{0{,}3}=2.
\]

\[
\boxed{2\ \kmh}
\]

\section{Производительность}

\subsection{Что такое производительность}

Производительность --- это скорость выполнения работы.

Например:
\[
6\ \text{вопросов в час}, \qquad
12\ \lpm, \qquad
5\ \text{деталей в минуту}.
\]

В задачах на производительность используется тот же принцип, что и в задачах на движение.

\begin{center}
\renewcommand{\arraystretch}{1.35}
\begin{tabularx}{0.86\linewidth}{X X}
\toprule
\textbf{Движение} & \textbf{Производительность} \\
\midrule
скорость $v$ & производительность $p$ \\
время $t$ & время $t$ \\
путь $S$ & результат $A$ \\
\bottomrule
\end{tabularx}
\end{center}

Главная формула:
\[
A=pt.
\]

Отсюда:
\[
p=\frac{A}{t}, \qquad t=\frac{A}{p}.
\]

\subsection{Таблица для задач на работу}

\begin{center}
\renewcommand{\arraystretch}{1.35}
\begin{tabularx}{0.9\linewidth}{X c c c}
\toprule
\textbf{Исполнитель} & \textbf{Производительность} & \textbf{Время} & \textbf{Результат} \\
\midrule
первый & $p_1$ & $t_1$ & $p_1t_1$ \\
второй & $p_2$ & $t_2$ & $p_2t_2$ \\
вместе & $p_1+p_2$ & $t$ & $(p_1+p_2)t$ \\
\bottomrule
\end{tabularx}
\end{center}

\important{Если исполнители работают вместе, складываются именно производительности.}

\subsection{Пример 1. Два мастера работают вместе}

Один мастер выполняет заказ за $3$ часа, а другой --- за $6$ часов. За сколько часов они выполнят заказ, работая вместе?

Результат во всех случаях --- один заказ.

\begin{center}
\renewcommand{\arraystretch}{1.35}
\begin{tabularx}{0.82\linewidth}{X c c c}
\toprule
 & $p$ & $t$ & $A$ \\
\midrule
первый мастер & $\dfrac{1}{3}$ & $3$ & $1$ \\
второй мастер & $\dfrac{1}{6}$ & $6$ & $1$ \\
вместе & $\dfrac{1}{x}$ & $x$ & $1$ \\
\bottomrule
\end{tabularx}
\end{center}

Общая производительность равна сумме производительностей:
\[
\frac{1}{x}=\frac{1}{3}+\frac{1}{6}.
\]

\[
\frac{1}{x}=\frac{2}{6}+\frac{1}{6}=\frac{3}{6}=\frac{1}{2}.
\]

\[
x=2.
\]

\[
\boxed{2\ \ch}
\]

\subsection{Пример 2. Количество вопросов в тесте}

Петя решает $6$ вопросов в час, а Ваня --- $7$ вопросов в час. Петя закончил тот же тест на $20$ минут позже Вани. Сколько вопросов было в тесте?

Пусть $x$ --- количество вопросов в тесте.

\begin{center}
\renewcommand{\arraystretch}{1.35}
\begin{tabularx}{0.82\linewidth}{X c c c}
\toprule
 & $p$ & $t$ & $A$ \\
\midrule
Петя & $6$ & $\dfrac{x}{6}$ & $x$ \\
Ваня & $7$ & $\dfrac{x}{7}$ & $x$ \\
\bottomrule
\end{tabularx}
\end{center}

Петя работает медленнее, значит, времени тратит больше.

Переводим:
\[
20\ \minut=\frac{20}{60}=\frac{1}{3}\ \ch.
\]

Составляем уравнение:
\[
\frac{x}{6}-\frac{x}{7}=\frac{1}{3}.
\]

\[
\frac{7x-6x}{42}=\frac{1}{3}.
\]

\[
\frac{x}{42}=\frac{1}{3}.
\]

\[
x=14.
\]

\[
\boxed{14\ \text{вопросов}}
\]

\section{Трубы, насосы и бассейны}

\subsection{Как понимать такие задачи}

В задачах с трубами и насосами производительность часто измеряется в литрах в минуту:
\[
\lpm.
\]

Если труба пропускает $12\ \lpm$, это означает:
\[
12\ \text{литров за } 1 \text{ минуту}.
\]

Если две трубы наполняют бассейн одновременно, их производительности складываются:
\[
p_{\text{общ}}=p_1+p_2.
\]

Если одна труба наполняет, а другая откачивает, производительности вычитаются:
\[
p_{\text{итог}}=p_{\text{наполнение}}-p_{\text{откачивание}}.
\]

\subsection{Пример 1. Одинаковое время работы}

Первая труба пропустила $770$ литров воды, а вторая за то же время пропустила $830$ литров. Вторая труба пропускает на $6\ \lpm$ больше первой. Сколько литров в минуту пропускает первая труба?

Пусть $x$ --- производительность первой трубы. Тогда производительность второй:
\[
x+6.
\]

\begin{center}
\renewcommand{\arraystretch}{1.35}
\begin{tabularx}{0.82\linewidth}{X c c c}
\toprule
 & $p$ & $t$ & $A$ \\
\midrule
первая труба & $x$ & $\dfrac{770}{x}$ & $770$ \\
вторая труба & $x+6$ & $\dfrac{830}{x+6}$ & $830$ \\
\bottomrule
\end{tabularx}
\end{center}

Время работы одинаковое:
\[
\frac{770}{x}=\frac{830}{x+6}.
\]

\[
770(x+6)=830x.
\]

\[
770x+4620=830x.
\]

\[
60x=4620.
\]

\[
x=77.
\]

\[
\boxed{77\ \lpm}
\]

\subsection{Пример 2. Совместная работа труб}

Первая труба наполняет бассейн за $21$ минуту, а вторая --- за $42$ минуты. За сколько минут они наполнят бассейн вместе?

Пусть весь бассейн --- это $1$.

\begin{center}
\renewcommand{\arraystretch}{1.35}
\begin{tabularx}{0.82\linewidth}{X c c c}
\toprule
 & $p$ & $t$ & $A$ \\
\midrule
первая труба & $\dfrac{1}{21}$ & $21$ & $1$ \\
вторая труба & $\dfrac{1}{42}$ & $42$ & $1$ \\
вместе & $\dfrac{1}{x}$ & $x$ & $1$ \\
\bottomrule
\end{tabularx}
\end{center}

\[
\frac{1}{x}=\frac{1}{21}+\frac{1}{42}.
\]

\[
\frac{1}{x}=\frac{2}{42}+\frac{1}{42}=\frac{3}{42}=\frac{1}{14}.
\]

\[
x=14.
\]

\[
\boxed{14\ \minut}
\]

\section{Алгоритм решения}

\begin{enumerate}
    \item Определите тип задачи: средняя скорость, поезд, догон протяжённых объектов, производительность.
    \item Составьте таблицу: скорость--время--путь или производительность--время--результат.
    \item Следите за единицами измерения.
    \item Если скорость дана в $\kmh$, время переводите в часы, а метры --- в километры.
    \item Для средней скорости используйте формулу:
    \[
    v_{\text{ср}}=\frac{S_{\text{общ}}}{t_{\text{общ}}}.
    \]
    \item Для протяжённых объектов учитывайте сумму длин.
    \item При движении навстречу скорости складывайте.
    \item При догоне скорости вычитайте.
    \item В задачах на производительность складывайте производительности, а не времена.
    \item Проверяйте, в каких единицах нужно дать ответ.
\end{enumerate}

\section{Типичные ошибки}

\begin{enumerate}
    \item \important{Считать среднюю скорость как среднее арифметическое всегда.} Это верно только для равных промежутков времени.
    \item \important{Перепутать половину времени и половину пути.} Эти случаи дают разные уравнения и разные ответы.
    \item \important{Забыть длину поезда или платформы.} Протяжённый объект должен пройти не точку, а сумму длин.
    \item \important{Сложить скорости при догоне.} При догоне используется разность скоростей.
    \item \important{Не перевести секунды в часы.} При скорости в $\kmh$ секунды нужно делить на $3600$.
    \item \important{Не перевести метры в километры.} Если скорость дана в $\kmh$, длину удобно переводить в километры.
    \item \important{Складывать времена при совместной работе.} При совместной работе складываются производительности.
    \item \important{Неверно выбрать объект поиска.} Часто лучше обозначать за $x$ именно то, что спрашивают в задаче.
\end{enumerate}

\newpage

\section{Тренировочный блок}

Решите задачи. Ответы приведены на отдельной странице после тренировочного блока.

\subsection*{Средний уровень}

\begin{enumerate}
    \item Первую половину времени автомобиль ехал со скоростью $72\ \kmh$, а вторую половину времени --- со скоростью $54\ \kmh$. Найдите среднюю скорость автомобиля.

    \item Первую половину пути автомобиль ехал со скоростью $60\ \kmh$, а вторую половину пути --- со скоростью $90\ \kmh$. Найдите среднюю скорость автомобиля.

    \item Поезд, скорость которого равна $72\ \kmh$, проходит мимо столба за $20$ секунд. Найдите длину поезда в метрах.
\end{enumerate}

\subsection*{Выше среднего уровня}

\begin{enumerate}[resume]
    \item Поезд, скорость которого равна $54\ \kmh$, проходит мимо платформы длиной $450\ \m$ за $1$ минуту. Найдите длину поезда в метрах.

    \item Поезд и пешеход движутся навстречу друг другу. Скорость поезда равна $72\ \kmh$, скорость пешехода равна $6\ \kmh$. Поезд проходит мимо пешехода за $30$ секунд. Найдите длину поезда в метрах.

    \item Два поезда длиной $400\ \m$ и $500\ \m$ движутся навстречу друг другу со скоростями $72\ \kmh$ и $90\ \kmh$. За сколько секунд они пройдут друг мимо друга?

    \item Короткая баржа длиной $40\ \m$ сначала отстаёт от длинной баржи длиной $60\ \m$ на $200\ \m$. Затем короткая баржа догоняет длинную и перегоняет её на $300\ \m$. На это уходит $18$ минут. На сколько $\kmh$ скорость короткой баржи больше скорости длинной?

    \item Один мастер выполняет заказ за $6$ часов, а другой --- за $12$ часов. За сколько часов они выполнят заказ, работая вместе?

    \item Петя решает $5$ вопросов в час, а Ваня --- $8$ вопросов в час. Петя закончил тот же тест на $45$ минут позже Вани. Сколько вопросов было в тесте?
\end{enumerate}

\subsection*{Сложный уровень}

\begin{enumerate}[resume]
    \item Первая труба пропустила $770$ литров воды, а вторая за то же время пропустила $830$ литров. Вторая труба пропускает на $6\ \lpm$ больше первой. Сколько литров в минуту пропускает первая труба?
\end{enumerate}

\newpage

\section{Ответы к тренировочному блоку}

\begin{center}
\renewcommand{\arraystretch}{1.35}
\begin{tabularx}{0.9\linewidth}{c X}
\toprule
\textbf{№} & \textbf{Ответ} \\
\midrule
1 & $63\ \kmh$ \\
2 & $72\ \kmh$ \\
3 & $400\ \m$ \\
4 & $450\ \m$ \\
5 & $650\ \m$ \\
6 & $20\ \ssec$ \\
7 & $2\ \kmh$ \\
8 & $4\ \ch$ \\
9 & $10$ вопросов \\
10 & $77\ \lpm$ \\
\bottomrule
\end{tabularx}
\end{center}

\end{document}